\section{Proofs}
\label{app:proofs}

The theorem and corollary statements remain in the main text. Their proofs are collected
here so that the main narrative can emphasize the representation chain, geometric
controls, and numerical evidence.

\subsection{Spectral projection and finite hopping support}
\label{app:proof-projection-range}

\begin{proof}[Proof of Proposition~\ref{prop:projection-destroys-range}]
Set $z=e^{ik}$. If $d(k)$ were a finite Fourier polynomial, there would be a Laurent
polynomial $f(z)$ whose restriction to the unit circle equals $d(k)$. Squaring on the
unit circle and using equality on infinitely many points would give the Laurent-polynomial
identity
\begin{equation}
 f(z)^2=\Delta^2+2t^2+t^2(z+z^{-1}).
 \label{eq:laurent-square}
\end{equation}
If $m$ and $n$ are the smallest and largest exponents with nonzero coefficients in
$f$, then the smallest and largest exponents in $f^2$ are $2m$ and $2n$; their
coefficients cannot cancel because each is the square of an extreme coefficient. The
right-hand side of Eq.~\eqref{eq:laurent-square} has extreme exponents $-1$ and $1$
when $t\ne0$, which is impossible because $2m$ and $2n$ are even. Thus $d$ is not a
finite Fourier polynomial. Equation~\eqref{eq:projected-range} has diagonal entries
containing $-d(k)/2$ plus finite Fourier terms, so the embedded projected operator is not
finite range.
\end{proof}

\subsection{Flat-parent residual--locality identity}
\label{app:proof-flat-parent}

\begin{proof}[Proof of Theorem~\ref{thm:flat-parent}]
A constant real rotation diagonalizes $D$ without changing any Frobenius norm, Fourier
weight, or minimization in the statement, so take $D=\delta\sigma_z$ with
$\delta\ne0$. Conjugation through a real rotation gives
\begin{equation}
 B_a^{(d)}(\mathbf q)
 =\delta\left[
 \cos\!\left(2\theta_a^{(d)}(\mathbf q)\right)\sigma_z
 -\sin\!\left(2\theta_a^{(d)}(\mathbf q)\right)\sigma_x
 \right].
 \label{eq:flat-rotated-pauli}
\end{equation}
Using the normalized integral representation
\begin{equation}
 J_0(x)=\int_{-1/2}^{1/2}e^{ix\sin(2\pi q)}\,dq
 \label{eq:bessel-integral}
\end{equation}
and the product structure of $\T^d$,
\begin{align}
 \int_{\T^d}e^{2i\theta_a^{(d)}(\mathbf q)}\,d\mathbf q
 &=\prod_{\alpha=1}^d
   \int_{-1/2}^{1/2}
   e^{i(2a/d)\sin(2\pi q_\alpha)}\,dq_\alpha\\
 &=J_0\!\left(\frac{2a}{d}\right)^d=c_d(a).
\end{align}
Global inversion $\mathbf q\mapsto-\mathbf q$ makes the sine average in
Eq.~\eqref{eq:flat-rotated-pauli} vanish, proving
Eq.~\eqref{eq:onsite-average}. Moreover,
$0\le2a/d\le\pi/4$, and the real part of the integrand in
Eq.~\eqref{eq:bessel-integral} is at least $\cos(2a/d)>0$; hence $c_d(a)>0$.

For any constant matrix $C$, the zero-mean property of
$B_a^{(d)}-B_{\mathbf0}^{(d)}$ gives the orthogonal decomposition
\begin{equation}
 \int_{\T^d}\norm{C-B_a^{(d)}(\mathbf q)}_F^2\,d\mathbf q
 =\int_{\T^d}\norm{B_a^{(d)}(\mathbf q)-B_{\mathbf0}^{(d)}}_F^2\,d\mathbf q
  +\norm{C-B_{\mathbf0}^{(d)}}_F^2.
 \label{eq:constant-projection}
\end{equation}
The constant $B_{\mathbf0}^{(d)}=c_d(a)D$ belongs to the allowed family in
Eq.~\eqref{eq:flat-operator-loss}: choose $s=0$, $\lambda=c_d(a)$, and $\phi=0$.
It is therefore the best constant restricted alignment. Equivalently, the normalized
maximum correlation with one global rotation is $c_d(a)$.

Parseval's identity applied entrywise yields
\begin{equation}
 \int_{\T^d}\norm{B_a^{(d)}-B_{\mathbf0}^{(d)}}_F^2
 =\sum_{\mathbf R\ne\mathbf0}\norm{B_{\mathbf R}^{(d)}}_F^2.
\end{equation}
Pointwise unitary invariance gives
$\norm{B_a^{(d)}(\mathbf q)}_F=\norm D_F$. A second use of Parseval, together with
Eq.~\eqref{eq:onsite-average}, therefore gives
\begin{equation}
 \sum_{\mathbf R\ne\mathbf0}\norm{B_{\mathbf R}^{(d)}}_F^2
 =\norm D_F^2-\norm{B_{\mathbf0}^{(d)}}_F^2
 =\norm D_F^2\left[1-c_d(a)^2\right].
\end{equation}
Dividing by $2E_{\mathrm{ref}}^2$ proves Eqs.~\eqref{eq:flat-residual} and
\eqref{eq:flat-tail}. Finally, fixing $s=0$ and $\lambda=1$ leaves the non-onsite
weight plus the onsite mismatch:
\begin{align}
 \inf_\phi L_{\mathrm{op},\phi}^2(0,1)
 &=\kappa[1-c_d(a)^2]+\kappa[1-c_d(a)]^2\\
 &=2\kappa[1-c_d(a)],
\end{align}
which proves Eq.~\eqref{eq:flat-exact-spectrum}.
\end{proof}

\subsection{Orthogonal shell coefficients and factorial tail}
\label{app:proof-orthogonal-tail}

\begin{proof}[Proof of Theorem~\ref{thm:orthogonal-fourier-tail}]
The Jacobi--Anger identity gives
\begin{equation}
 e^{2i\theta_a^{(d)}(\mathbf q)}
 =\prod_{\alpha=1}^d
  \sum_{n\in\mathbb Z}J_n(z)e^{2\pi inq_\alpha}
 =\sum_{\mathbf R\in\mathbb Z^d}
  j_{\mathbf R}(z)e^{2\pi i\mathbf q\cdot\mathbf R}.
\end{equation}
Because $J_{-n}(z)=(-1)^nJ_n(z)$, the cosine coefficient vanishes for odd
$p(\mathbf R)$ and equals $j_{\mathbf R}(z)$ for even $p(\mathbf R)$; the sine
coefficient vanishes for even parity and equals $-ij_{\mathbf R}(z)$ for odd parity.
Substitution into Eq.~\eqref{eq:flat-rotated-pauli} proves
Eq.~\eqref{eq:orthogonal-fourier-block}. Since both Pauli matrices have squared
Frobenius norm two,
\begin{equation}
 \norm{B_{\mathbf R}^{(d)}}_F^2
 =\norm D_F^2\prod_{\alpha=1}^dJ_{R_\alpha}(z)^2.
 \label{eq:block-power-product}
\end{equation}
Equation~\eqref{eq:shell-power} follows by grouping a complete shell. Parseval applied
to the unit-modulus function $e^{iz\sin(2\pi q)}$ gives
$\sum_{n\in\mathbb Z}J_n(z)^2=1$; taking the $d$-fold product proves
Eq.~\eqref{eq:shell-partition}.

If $\norm{\mathbf R}_2^2>S$, at least one coordinate obeys $|R_\alpha|\ge L$.
The union bound and the one-dimensional Bessel partition therefore give
\begin{equation}
 \sum_{\norm{\mathbf R}_2^2>S}
 \prod_{\alpha=1}^dJ_{R_\alpha}(z)^2
 \le 2d\sum_{n=L}^{\infty}|J_n(z)|^2.
 \label{eq:tail-union}
\end{equation}
The defining Bessel series and $(n+m)!\ge n!m!$ imply, for $n\ge0$,
\begin{equation}
 |J_n(z)|
 \le \frac{x^n}{n!}\sum_{m=0}^{\infty}\frac{x^{2m}}{m!}
 =\frac{x^n}{n!}e^{x^2}.
\end{equation}
Moreover, $(L+m)!\ge L!(L+1)^m$, so when $x<L+1$,
\begin{equation}
 \sum_{n=L}^{\infty}\frac{x^{2n}}{(n!)^2}
 \le \frac{x^{2L}}{(L!)^2}
 \sum_{m=0}^{\infty}\left(\frac{x^2}{(L+1)^2}\right)^m
 =\frac{x^{2L}}
 {(L!)^2\left[1-x^2/(L+1)^2\right]}.
\end{equation}
Combining this estimate with Eq.~\eqref{eq:tail-union}, using
$e^{2x^2}=e^{z^2/2}$ and Eq.~\eqref{eq:block-power-product}, proves
Eq.~\eqref{eq:orthogonal-factorial-tail}.
\end{proof}

\subsection{Monotonicity and small-amplitude dimension law}
\label{app:proof-dimension-law}

\begin{proof}[Proof of Corollary~\ref{cor:dimension-law}]
For $x>0$, $J_0'(x)=-J_1(x)$. On $0<x\le\pi/4$, the alternating series
\begin{equation}
 J_1(x)=\sum_{m=0}^{\infty}
 \frac{(-1)^m}{m!(m+1)!}\left(\frac{x}{2}\right)^{2m+1}
\end{equation}
has strictly decreasing term magnitudes because their successive ratio is at most
$\pi^2/128<1$. Hence $J_1(x)>0$. Together with the positivity already proved for
$J_0(2a/d)$, this shows
\begin{equation}
 c_d'(a)=-2J_0(2a/d)^{d-1}J_1(2a/d)<0
\end{equation}
for $a>0$. The threshold monotonicity follows immediately. Finally,
$J_0(x)=1-x^2/4+x^4/64+O(x^6)$. Substituting $x=2a/d$, raising to the $d$th power, and
expanding the two threshold formulas proves
Eqs.~\eqref{eq:correlation-small-a}--\eqref{eq:exact-threshold-small-a}.
\end{proof}

\subsection{Bond-orbit flat-parent comparator}
\label{app:proof-bond-orbit}

\begin{proof}[Proof of Theorem~\ref{thm:bond-orbit}]
Equation~\eqref{eq:bond-orbit-phase} is odd under
$\mathbf q\mapsto-\mathbf q$, so the imaginary part of
Eq.~\eqref{eq:bond-orbit-correlation} integrates to zero. Since
$|2\theta_{a,\mathcal B}|\le\pi/2$, its cosine is nonnegative; it is positive in a
neighborhood of $\mathbf q=\mathbf0$, where the phase vanishes. Hence the integral is
strictly positive. Diagonalizing $D$ by one constant real rotation and repeating
Eq.~\eqref{eq:flat-rotated-pauli} gives
$D_{\mathbf0}^{\mathcal B}=c_{\mathcal B,\boldsymbol\eta}D$. The constant-projection
identity in Eq.~\eqref{eq:constant-projection} and entrywise Parseval then prove the two
residual formulas exactly as in Theorem~\ref{thm:flat-parent}.

For the explicit comparator, apply the Jacobi--Anger expansion to every bond harmonic:
\begin{equation}
 e^{i(2a\eta_j/W)\sin(2\pi\mathbf q\cdot\mathbf b_j)}
 =\sum_{n_j\in\Z}
 J_{n_j}\!\left(\frac{2a\eta_j}{W}\right)
 e^{2\pi i n_j\mathbf q\cdot\mathbf b_j}.
\end{equation}
The finite product of these absolutely convergent series may be integrated termwise.
Torus orthogonality retains exactly the tuples satisfying
$\sum_jn_j\mathbf b_j=\mathbf0$, which proves
Eq.~\eqref{eq:bond-orbit-bessel-sum}. If the integer kernel of the bond matrix is
trivial, only $\mathbf n=\mathbf0$ remains, giving
Eq.~\eqref{eq:independent-bond-product}.
\end{proof}

\subsection{Constrained-series complete-shell certificate}
\label{app:proof-constrained-tail}

\begin{proof}[Proof of Theorem~\ref{thm:constrained-tail}]
Jacobi--Anger expansion and torus orthogonality give
Eq.~\eqref{eq:constrained-coefficient}. The symmetry
$e^{2i\theta(-\mathbf q)}=\overline{e^{2i\theta(\mathbf q)}}$ makes every
$g_{\mathbf R}$ real. Splitting the cosine and sine coefficients into
$(g_{\mathbf R}\pm g_{-\mathbf R})/2$ proves
Eq.~\eqref{eq:constrained-matrix-block}. Hence
\begin{equation}
 \frac{\norm{D_{\mathbf R}^{\mathcal B}}_F^2}{\norm D_F^2}
 =\frac{g_{\mathbf R}^2+g_{-\mathbf R}^2}{2}.
\end{equation}
Parseval for the unit-modulus scalar phase gives
$\sum_{\mathbf R}g_{\mathbf R}^2=1$. Because $\mathcal S_S$ is inversion complete,
the exact matrix tail normalized by $\norm D_F^2$ is
$1-\sum_{\mathbf R\in\mathcal S_S}g_{\mathbf R}^2$. The reverse triangle inequality
$|g_{\mathbf R}|\ge\max\{0,|\widetilde g_{\mathbf R}|-\epsilon_{\mathbf R}\}$
proves Eq.~\eqref{eq:constrained-tail-certificate}.

Solving the triangular constraint gives
$(n_1,n_2,n_3)=(r_1-\ell,r_2+\ell,\ell)$, which proves
Eq.~\eqref{eq:triangular-coefficient}. Taking
$(u,v,w)=(n_4,n_5,n_6)$ in the FCC constraint gives
\begin{equation}
 (n_1,n_2,n_3)=(r_1+u+v,r_2-u+w,r_3-v-w),
\end{equation}
which proves Eq.~\eqref{eq:fcc-coefficient}. Finally, the defining Bessel series gives
$|J_n(z)|\le e^{x^2}x^{|n|}/|n|!$. Summing it over all integers gives $A_0$; summing
over $|n|>K$ and bounding the remaining exponential-series tail geometrically gives
$A_K$. In the triangular remainder the other two Bessel factors have magnitude at most
one. In the FCC remainder, a union bound over the three free indices gives
$3A_KA_0^2$.
\end{proof}


\section{Diagnostic definitions}
\label{app:diagnostics}

Every discrepancy below is defined only after its two arguments have been represented on
the same sampled reciprocal mesh, in the same retained state space, basis, gauge, energy
reference, units, and lattice geometry. A small numerical value cannot repair a missing
state-space identification.

\subsection{Pointwise matrix and sorted-band errors}

Let $A_\ell,B_\ell\in\mathbb C^{r\times r}$ be Hermitian fields evaluated at the same
$M$ reciprocal points, with ordered eigenvalues
$\varepsilon_j^A(\mathbf q_\ell)$ and
$\varepsilon_j^B(\mathbf q_\ell)$. Define the rank-normalized pointwise Hamiltonian
defect
\begin{equation}
 d_{H,\ell}(A,B)=\frac{\norm{A_\ell-B_\ell}_F}{\sqrt r}.
 \label{eq:app-pointwise-hamiltonian-defect}
\end{equation}
The four primary error channels are
\begin{align}
 \epsilon_{H,\max}(A,B)&=\max_\ell d_{H,\ell}(A,B),\\
 \epsilon_{H,\mathrm{rms}}(A,B)
 &=\left[\frac1M\sum_{\ell=1}^M d_{H,\ell}(A,B)^2\right]^{1/2},\\
 \epsilon_{b,\max}(A,B)
 &=\max_{\ell,j}\left|\varepsilon_j^A(\mathbf q_\ell)
                   -\varepsilon_j^B(\mathbf q_\ell)\right|,\\
 \epsilon_{b,\mathrm{rms}}(A,B)
 &=\left[\frac1{Mr}\sum_{\ell=1}^M\sum_{j=1}^r
  \left|\varepsilon_j^A(\mathbf q_\ell)
       -\varepsilon_j^B(\mathbf q_\ell)\right|^2\right]^{1/2}.
 \label{eq:app-primary-error-channels}
\end{align}
They have energy units. A displayed dimensionless value divides them by the declared
$E_{\mathrm{ref}}$. The present retained rank is $r=2$, so the implementation uses the
factor $1/\sqrt2$ in Eq.~\eqref{eq:app-pointwise-hamiltonian-defect}. Pointwise Weyl
control is checked separately; in this normalization,
\begin{equation}
 \epsilon_{b,\max}(A,B)\le
 \sqrt r\,\epsilon_{H,\max}(A,B).
 \label{eq:app-weyl-check}
\end{equation}
No equality between matrix and band errors is assumed.

\subsection{Interpolation, deliberate truncation, and total error}

Write $H^{\mathrm{dir}}$ for the direct withheld-mesh field,
$H_N^{\mathrm{ord}}$ for the full ordinary Fourier interpolant, and
$H_{N,R}^{\mathrm{ord}}$ for its complete-radius truncation. For each of the four
functionals in Eq.~\eqref{eq:app-primary-error-channels}, the three approximation
channels are
\begin{align}
 \mathcal E_{\mathrm{interp}}&=\mathcal E(H_N^{\mathrm{ord}},H^{\mathrm{dir}}),\\
 \mathcal E_{\mathrm{trunc}}&=\mathcal E(H_{N,R}^{\mathrm{ord}},H_N^{\mathrm{ord}}),\\
 \mathcal E_{\mathrm{total}}&=\mathcal E(H_{N,R}^{\mathrm{ord}},H^{\mathrm{dir}}).
 \label{eq:three-error-channels}
\end{align}
These channels are not added: interpolation and truncation may reinforce or cancel. For
a strain-responsive control, the same definitions are applied to the gauge-aligned
difference field $\Delta H_s$ after constructing it in one common retained space. The
geometry-only strain control reports locality under a changing metric and is not a
three-channel strain-potential calculation.

\subsection{Fourier power, complete shells, and first passing}

For ordinary hopping blocks $T_{\mathbf n}$ and a complete physical-metric support
$\mathcal S_s$, define
\begin{align}
 P_{\mathrm{shell}}(s)
 &=\frac1r\sum_{\mathbf n\in\mathcal S_s\setminus\mathcal S_{s-1}}
   \norm{T_{\mathbf n}}_F^2,\\
 P_{\mathrm{omit}}(s)
 &=\frac1r\sum_{\mathbf n\notin\mathcal S_s}
   \norm{T_{\mathbf n}}_F^2.
 \label{eq:app-shell-power}
\end{align}
Division by $E_{\mathrm{ref}}^2$ gives the dimensionless power. Cyclic Nyquist power is
reported separately and is not silently reassigned to an ordinary conjugate-closed
hopping block. Parseval defects compare reciprocal-space power with the complete cyclic
block inventory; truncation Parseval defects compare deliberate omitted power with its
reciprocal-space reconstruction consequence.

For diagnostic channel $X$ and tolerance $\tau$, the sampled first-passing shell is
\begin{equation}
 \widehat s_X(\tau)=
 \min\{s:\epsilon_{X,\mathrm{total},\max}(s)\le\tau\}.
 \label{eq:app-first-passing}
\end{equation}
The record retains the physical radius, block count, passing error, and immediately
preceding error. If no complete shell passes before the declared cap, the decision is
unresolved rather than assigned to the cap.

\subsection{Alignment, reconstruction, and isolation diagnostics}

The reconstruction ledger keeps parent Hermiticity, complete-eigensystem reconstruction,
frame orthonormality, projector reconstruction, parent--retained lifting, Fourier
reconstruction, covariant equality, and restricted global alignment as separate
defects. External retained-subspace isolation is reported both as the minimum sampled
gap and as a continuous lower bound. A sampled gap is not relabeled as a continuous
certificate.

Restricted alignment records one global rotation, the optimized nonnegative splitting
scale, the finite-mesh maximum correlation, the minimum operator threshold, and the
distinct unit-splitting exact-spectrum joint threshold. A fixed-unit-scale rotation fit
cannot substitute for the optimized-scale threshold. Flat-family analytic correlations,
quadrature defects, residual--tail identities, amplitude midpoint defects, and
amplitude-resolution ladders remain separate checks.

\subsection{Strain candidates, sampled envelopes, and transition brackets}

At sampled strain $s_i$, let $\kappa_i$ and $\mu_i$ denote the restricted-alignment
splitting scale and maximum correlation. The pointwise optimized requirement and the
fixed reference-candidate requirement are
\begin{align}
 \tau_{o,\min}(s_i)
 &=\sqrt{\kappa_i[1-\max(0,\mu_i)^2]},\\
 \tau_{o,p_0}(s_i)
 &=\sqrt{2\kappa_i(1-\mu_i)},
 \qquad p_0=(0,1).
 \label{eq:app-strain-candidate-requirements}
\end{align}
The first quantity permits an independently optimized candidate at every sampled
strain. It is not one common witness. The second evaluates one prescribed parameter
path, but it is not a witness until declared spectral and operator limits establish
feasibility. Their hatted maxima, and the hatted locality radius and block-count maxima,
are finite-sample statistics:
\begin{equation}
 \widehat\tau_{o,\min}=\max_i\tau_{o,\min}(s_i),\quad
 \widehat\tau_{o,p_0}=\max_i\tau_{o,p_0}(s_i),\quad
 \widehat R_\varepsilon=\max_iR_\varepsilon(s_i),\quad
 \widehat N_\varepsilon=\max_iN_\varepsilon(s_i).
 \label{eq:app-strain-sampled-envelopes}
\end{equation}
Each maximum retains its own sample location; the locations need not agree.

The response-sweep qualification envelope additionally retains
\begin{equation}
 \min_i g_{\mathrm{ext}}(s_i),\qquad
 \min_i\sigma_{\mathrm{seed}}(s_i),\qquad
 \min_i\sigma_{\mathrm{frame}}(s_i),\qquad
 \max_i r_{\mathrm{point}}(s_i),\qquad
 \max_i r_{\mathrm{global}}(s_i),
 \label{eq:app-strain-qualification-envelope}
\end{equation}
where $g_{\mathrm{ext}}$ is the sampled external isolation gap,
$\sigma_{\mathrm{seed}}$ is the minimum seed-projector overlap,
$\sigma_{\mathrm{frame}}$ is the minimum singular value of the pointwise frame overlap,
and the two $r$ channels are pointwise and restricted-global frame-alignment residuals.
Every extremum retains the first sampled strain at which it occurs. These quantities
qualify the retained state-space identification; they are not response observables.

For the frozen-parent and radial-response controls, the sweep also takes separate sampled
maxima of all four errors in Eq.~\eqref{eq:app-primary-error-channels} for the aligned
response, interpolation, deliberate truncation, and total channels. The same reduction
is applied to $\Delta H_s$, ordinary omitted power, cyclic Nyquist power, scalar and
traceless response components, onsite and offsite response power, and minimum and
maximum prescribed hopping-response factors. No maxima from distinct channels are
combined into a synthetic sample, and none is relabeled as a continuum supremum.

A detected transition bracket is an adjacent sampled interval whose endpoint values of
a discrete channel differ. The present discrete channels are first-passing block count
at each tail tolerance and number of distinct bond-length orbits. Five bisections reduce
a base-grid interval of width at most $1/200$ to width at most $1/6400$. The bracket
localizes a detected change only. Equal endpoint values do not prove that the channel is
constant inside the interval, and the bracket supplies no derivative or interval bound
for the continuous path.

\subsection{Transport and closure diagnostics}

For rank-two closure diagnostics, let $W_\alpha(\mathbf k_\perp)$ be the ordered product
of full-rank transport links around reciprocal direction $\alpha$ at transverse base
point $\mathbf k_\perp$. If
$-\pi\le\phi_{\alpha,0}(\mathbf k_\perp)\le
\phi_{\alpha,1}(\mathbf k_\perp)\le\pi$ are its ordered principal eigenphases, define
\begin{equation}
 \Phi_j=\max_{\alpha,\mathbf k_\perp}
 |\phi_{\alpha,j}(\mathbf k_\perp)|,
 \qquad j\in\{0,1\}.
 \label{eq:multidimensional-closure-phase-reduction}
\end{equation}
This reduction is applicable only when every contributing transport link is full rank.
It is accompanied by the minimum transport singular value, Procrustes residual,
rank-cutoff and roundoff controls, link-unitarity defect, closure-unitarity defect, and
maximum plaquette phase. These closure and plaquette channels have no asserted
topological meaning here. A rank-deficient link makes the dependent closure channels
explicitly inapplicable rather than numerically zero.

\subsection{Numerical validation boundary}

The numerical checks also retain mesh separation, metric covariance, basis-distance
extrema where applicable, refinement changes, and independent reconstruction defects.
Agreement of these diagnostics verifies the declared finite calculation. It does not
establish a continuum limit, supply an externally chosen adequacy tolerance, or convert
a synthetic model into a material-specific Hamiltonian.


\section{Coordinate, Fourier, and state-space conventions}
\label{app:conventions}

Let the columns of $A\in\mathbb R^{d\times d}$ be direct primitive vectors. Reduced
reciprocal coordinates $\mathbf q\in\T^d$ and integer lattice modes
$\mathbf n\in\Z^d$ are mapped to Cartesian vectors by
\begin{equation}
 B=2\pi A^{-\mathsf T},\qquad
 \mathbf k=B\mathbf q,\qquad
 \mathbf R=A\mathbf n.
 \label{eq:app-direct-reciprocal-map}
\end{equation}
Consequently $\mathbf k\cdot\mathbf R=2\pi\mathbf q\cdot\mathbf n$. The
nonorthogonal protocols retain this reduced-coordinate convention; the physical radius
of a hopping mode is $\norm{A\mathbf n}$ rather than a coordinate-axis norm.

For a matrix field $H(\mathbf q)$, the continuous Fourier convention is
\begin{equation}
 H_{\mathbf n}=\int_{\T^d}
 e^{-2\pi i\mathbf q\cdot\mathbf n}H(\mathbf q)\,d\mathbf q,
 \qquad
 H(\mathbf q)=\sum_{\mathbf n\in\Z^d}
 e^{2\pi i\mathbf q\cdot\mathbf n}H_{\mathbf n}.
 \label{eq:app-fourier-convention}
\end{equation}
A finite even mesh first produces a cyclic centered inventory. Cyclic Nyquist modes are
reported separately before selecting an ordinary conjugate-closed hopping inventory.
Every finite support is completed in the physical metric and retains all equal-radius
ties.

Hamiltonian subtraction is performed only after state spaces have been identified. The
parent orbital field, complete eigensystem, retained projector, periodic retained frame,
and Wannier blocks remain distinct representations. Pointwise equality of sorted
eigenvalues is not used as a substitute for an operator identification.

\section{Model inventory and numerical design}
\label{app:protocols}

The fixed-geometry comparison contains four model families, and the strain study adds
two geometry paths.

\begin{table}[htbp]
\centering
\small
\caption{Model and control inventory.}
\label{tab:app-protocol-inventory}
\begin{tabular}{@{}>{\raggedright\arraybackslash}p{0.30\textwidth}
                    >{\raggedright\arraybackslash}p{0.62\textwidth}@{}}
\toprule
Family or path & Geometry and role \\
\midrule
Square 2-D & Orthogonal directed-parent baseline \\
Triangular/honeycomb 2-D & Nonorthogonal target with matched square and triangular
Bravais controls \\
Simple-cubic 3-D & Orthogonal directed-parent baseline \\
FCC/diamond 3-D & Nonorthogonal target with matched simple-cubic and FCC Bravais
controls \\
Rhombic strain & Unit-area triangular-to-square path with three controls \\
Bain strain & Unit-volume FCC-to-BCC path with three controls \\
\bottomrule
\end{tabular}
\end{table}

The construction and shifted evaluation meshes are tensor products. The evaluation
shift is one third of one evaluation-grid interval in every reduced reciprocal
coordinate. Table~\ref{tab:app-numerical-design} gives the three-level mesh ladder. The
reported $10^{-3}$ and $10^{-5}$ tail-error tolerances select first-passing complete
supports; they are support-selection coordinates, not acceptance limits for restricted
operator mismatch.

\begin{table}[htbp]
\centering
\small
\caption{Fixed-geometry numerical design. Mesh sizes are per reciprocal coordinate.}
\label{tab:app-numerical-design}
\begin{tabular}{@{}lccc@{}}
\toprule
Family & Construction sizes & Evaluation sizes & Shell cap \\
\midrule
Square 2-D & $16,32,64$ & $33,65,129$ & 25 \\
Triangular/honeycomb 2-D & $16,32,64$ & $33,65,129$ & 21 \\
Simple-cubic 3-D & $8,16,32$ & $17,33,65$ & 9 \\
FCC/diamond 3-D & $8,16,32$ & $17,33,65$ & 7 \\
\bottomrule
\end{tabular}
\end{table}

The orthogonal gauge amplitude is sampled at $0$, $\pi/8$, and $\pi/4$. The
triangular/honeycomb family uses maximum amplitude $\pi/(4\sqrt6)$ and the FCC/diamond
family uses $\pi/(4\sqrt{12})$; both include zero and the midpoint. Every family samples
parent scales $\rho\in\{0,1\}$. These are synthetic controls, not graphene or silicon
parameterizations.

\section{Cross-protocol requirement coordinates}
\label{app:tolerance-requirements}

The positive-splitting-scale optimization and the exact-spectrum unit-scale question are
kept separate. The earlier two-coordinate requirement reduction records
\begin{equation}
 x=\frac{\tau_{o,\min}}{\sqrt\kappa},\qquad
 y=\frac{\tau_{o,\mathrm{joint}}}{E_{\mathrm{ref}}}.
 \label{eq:app-requirement-pair}
\end{equation}
The componentwise maxima over cases, parent scales, and amplitudes need not originate
from the same sampled point and are not attainable decisions in general.
Table~\ref{tab:app-case-requirements} retains these scalar audit coordinates for each
case, but the main text replaces their visual envelope with the shell--scale Pareto
formulation in Eq.~\eqref{eq:pareto-decision-objectives}.

\begin{table}[htbp]
\centering
\small
\caption{Componentwise case envelopes for the sampled fixed-geometry calculation.
These coordinates summarize calculated requirements; they are not acceptance limits.}
\label{tab:app-case-requirements}
\begin{tabular}{@{}lcc@{}}
\toprule
Case & $\max\tau_{o,\min}/\sqrt\kappa$ &
$\max\tau_{o,\mathrm{joint}}/E_{\mathrm{ref}}$ \\
\midrule
Square directed parent & 0.688457 & 0.328854 \\
Matched square Bravais control & 0.586917 & 0.273379 \\
Triangular Bravais control & 0.584676 & 0.263702 \\
Honeycomb-inspired target & 0.468419 & 0.083542 \\
Simple-cubic directed parent & 0.584781 & 0.270259 \\
Matched simple-cubic Bravais control & 0.432793 & 0.195029 \\
FCC Bravais control & 0.430770 & 0.189573 \\
Diamond-inspired target & 0.347324 & 0.052890 \\
\bottomrule
\end{tabular}
\end{table}

Table~\ref{tab:app-case-requirements} reports requirement coordinates without deriving
acceptance limits from the observed values. Operating-point selection requires an
external application criterion.

\section{Case-resolved shell--scale decision landscapes}
\label{app:case-decision-landscapes}

Table~\ref{tab:app-decision-landscape-inventory} records the fixed-geometry endpoint
sampling used in the case-resolved plots. Each case uses its maximum declared parent
scale and gauge amplitude, every complete shell under its own physical metric and shell
cap, and 17 uniformly spaced values of the scale-path coordinate in
Eq.~\eqref{eq:scale-path}. Exact componentwise dominance is evaluated only among these
calculated vertices. The number of nondominated vertices is therefore a sampling
diagnostic, not a continuum invariant or a measure of scientific adequacy.

\begin{table}[htbp]
\centering
\small
\caption{Sampled shell--scale landscape inventory.}
\label{tab:app-decision-landscape-inventory}
\begin{tabular}{@{}lrrr@{}}
\toprule
Case & $\lambda_*$ & Calculated & Nondominated \\
\midrule
Square directed parent & 0.729890 & 238 & 120 \\
Matched square Bravais control & 0.813169 & 221 & 127 \\
Triangular Bravais control & 0.835297 & 187 & 156 \\
Graphene-inspired honeycomb & 0.883506 & 187 & 126 \\
Simple-cubic directed parent & 0.812597 & 153 & 85 \\
Matched simple-cubic Bravais control & 0.902347 & 119 & 80 \\
FCC Bravais control & 0.913939 & 136 & 116 \\
Silicon-inspired diamond & 0.937745 & 136 & 34 \\
\bottomrule
\end{tabular}
\end{table}

The parent scale is $1$ in every row. The gauge amplitude and complete-shell inventory
are as follows; parenthesized block counts give $N(S)$ in the same order as the listed
shell indices.
\begin{description}
 \item[Square directed parent:] $a=\pi/4$;
 $S=(0,\allowbreak 1,\allowbreak 2,\allowbreak 4,\allowbreak 5,\allowbreak 8,\allowbreak 9,\allowbreak 10,\allowbreak 13,\allowbreak 16,\allowbreak 17,\allowbreak 18,\allowbreak 20,\allowbreak 25)$;
 $N=(1,\allowbreak 5,\allowbreak 9,\allowbreak 13,\allowbreak 21,\allowbreak 25,\allowbreak 29,\allowbreak 37,\allowbreak 45,\allowbreak 49,\allowbreak 57,\allowbreak 61,\allowbreak 69,\allowbreak 81)$.
 \item[Matched square control:] $a=\pi/(4\sqrt6)$;
 $S=(0,\allowbreak 1,\allowbreak 2,\allowbreak 4,\allowbreak 5,\allowbreak 8,\allowbreak 9,\allowbreak 10,\allowbreak 13,\allowbreak 16,\allowbreak 17,\allowbreak 18,\allowbreak 20)$;
 $N=(1,\allowbreak 5,\allowbreak 9,\allowbreak 13,\allowbreak 21,\allowbreak 25,\allowbreak 29,\allowbreak 37,\allowbreak 45,\allowbreak 49,\allowbreak 57,\allowbreak 61,\allowbreak 69)$.
 \item[Triangular control:] $a=\pi/(4\sqrt6)$;
 $S=(0,\allowbreak 1,\allowbreak 3,\allowbreak 4,\allowbreak 7,\allowbreak 9,\allowbreak 12,\allowbreak 13,\allowbreak 16,\allowbreak 19,\allowbreak 21)$;
 $N=(1,\allowbreak 7,\allowbreak 13,\allowbreak 19,\allowbreak 31,\allowbreak 37,\allowbreak 43,\allowbreak 55,\allowbreak 61,\allowbreak 73,\allowbreak 85)$.
 \item[Honeycomb:] $a=\pi/(4\sqrt6)$; it uses the same shell indices and block counts
 as the triangular control.
 \item[Cubic parent:] $a=\pi/4$;
 $S=(0,1,2,3,4,5,6,8,9)$; $N=(1,7,19,27,33,57,81,93,123)$.
 \item[Matched cubic control:] $a=\pi/(4\sqrt{12})$;
 $S=(0,1,2,3,4,5,6)$; $N=(1,7,19,27,33,57,81)$.
 \item[FCC control:] $a=\pi/(4\sqrt{12})$;
 $S=(0,1,2,3,4,5,6,7)$; $N=(1,13,19,43,55,79,87,135)$.
 \item[Diamond:] $a=\pi/(4\sqrt{12})$; it uses the same shell indices and block
 counts as the FCC control.
\end{description}

The axes of every plot are the three components of
$\mathcal J_s(S,u)$ in Eq.~\eqref{eq:pareto-decision-objectives}, evaluated here at the
fixed-geometry endpoint amplitude rather than along a strain path. Color gives the
scale-path coordinate $u$. The displayed triangular faces connect only adjacent cells in
the sampled $(S,u)$ grid. They do not assert that the objective-space image is convex,
smooth, continuously covered, or valid between vertices.

\begin{figure}[p]
 \centering
 \includegraphics[width=0.92\textwidth]{figures/square-directed-parent-shell-scale-decision-landscape.pdf}
 \caption{Square directed-parent decision landscape. The comparatively broad separation
 between $\lambda_*$ and unity produces a visible operator--spectrum scale compromise.}
 \label{fig:app-decision-square-parent}
\end{figure}

\begin{figure}[p]
 \centering
 \includegraphics[width=0.92\textwidth]{figures/matched-square-bravais-control-shell-scale-decision-landscape.pdf}
 \caption{Matched square Bravais-control decision landscape. This control changes the
 parent law while retaining the square metric.}
 \label{fig:app-decision-square-control}
\end{figure}

\begin{figure}[p]
 \centering
 \includegraphics[width=0.92\textwidth]{figures/triangular-bravais-control-shell-scale-decision-landscape.pdf}
 \caption{Triangular Bravais-control decision landscape. Its shell inventory is defined
 by the non-orthogonal physical metric rather than square-coordinate radius.}
 \label{fig:app-decision-triangular-control}
\end{figure}

\begin{figure}[p]
 \centering
 \includegraphics[width=0.92\textwidth]{figures/graphene-inspired-honeycomb-shell-scale-decision-landscape.pdf}
 \caption{Graphene-inspired honeycomb decision landscape. ``Inspired'' denotes the
 synthetic geometry and connectivity, not a graphene material parametrization.}
 \label{fig:app-decision-honeycomb}
\end{figure}

\begin{figure}[p]
 \centering
 \includegraphics[width=0.92\textwidth]{figures/simple-cubic-directed-parent-shell-scale-decision-landscape.pdf}
 \caption{Simple-cubic directed-parent decision landscape.}
 \label{fig:app-decision-cubic-parent}
\end{figure}

\begin{figure}[p]
 \centering
 \includegraphics[width=0.92\textwidth]{figures/matched-simple-cubic-bravais-control-shell-scale-decision-landscape.pdf}
 \caption{Matched simple-cubic Bravais-control decision landscape.}
 \label{fig:app-decision-cubic-control}
\end{figure}

\begin{figure}[p]
 \centering
 \includegraphics[width=0.92\textwidth]{figures/fcc-bravais-control-shell-scale-decision-landscape.pdf}
 \caption{FCC Bravais-control decision landscape using complete FCC metric shells.}
 \label{fig:app-decision-fcc-control}
\end{figure}

\begin{figure}[p]
 \centering
 \includegraphics[width=0.92\textwidth]{figures/silicon-inspired-diamond-shell-scale-decision-landscape.pdf}
 \caption{Silicon-inspired diamond decision landscape. Its optimized scale is closest to
 unity and its sampled nondominated set is the most selective of the eight cases. This
 remains a synthetic rank-two control, not a silicon material calculation.}
 \label{fig:app-decision-diamond}
\end{figure}

\clearpage
\section{Finite-strain control definitions}
\label{app:strain-controls}

The strain sequence separates three distinct questions. It does not assign the geometry
control, null response, and prescribed response to a single model.

\subsection{Geometry-only control}

The 2-D unit-area rhombic path is
\begin{equation}
 A_2(\theta)=\frac{1}{\sqrt{\sin\theta}}
 \begin{pmatrix}1&\cos\theta\\0&\sin\theta\end{pmatrix},
 \qquad \theta:\pi/3\longrightarrow\pi/2.
 \label{eq:app-rhombic-path}
\end{equation}
The 3-D path is the unit-primitive-volume Bain deformation from FCC to BCC in a common
body-centered tetragonal primitive basis. In both cases, the reduced-coordinate parent,
imposed gauge, and Fourier-block weights remain fixed. The changing physical metric
alone reorders complete shells and changes first-passing physical radius and retained
block count.

\subsection{Frozen-parent null-response control}

The frozen-parent control uses the four-state bipartite parent and the
same imposed gauge at every strain point while keeping onsite and hopping matrices fixed.
Its gauge-aligned strain difference is expected to vanish up to numerical roundoff. This
is a null-response control, not a strain-dependent Hamiltonian model.

\subsection{Synthetic radial-response control}

The radial-response control changes only the selected inter-sublattice bond amplitudes:
\begin{equation}
 \frac{t_j(s)}{t_j(0)}=
 \exp\!\left[-\left(\frac{d_j(s)}{d_j(0)}-1\right)\right].
 \label{eq:app-radial-response}
\end{equation}
No onsite response, refit, structural relaxation, or material calibration is included.
The response field is aligned in a common retained space before scalar/traceless,
onsite/offsite, interpolation, deliberate-truncation, and total channels are reported.
The interpolation, truncation, and total channels are not added.

At the final 2-D metric, the calculated response factors span
$0.81758$--$1.27146$ and the maximum aligned retained-space perturbation is
$0.063575E_{\mathrm{ref}}$. At the BCC endpoint, the four selected bonds share response
$0.97564$ and the maximum perturbation is $0.007289E_{\mathrm{ref}}$. These are
results for the prescribed synthetic law, not material parameters or predictions.

\section{Independent numerical reconstruction and scope}
\label{app:verification}

An independent numerical oracle reconstructs mesh coordinates, minimum
construction--evaluation separation, parent and retained fields, restricted global
alignment, direct finite-sum Fourier blocks, Parseval identities, Hermiticity defects,
complete physical shells, interpolation/truncation/total matrix and band channels, Weyl
bounds, and first-passing shell decisions. It does not invoke the production Fourier
implementation it verifies. Nonorthogonal cases retain reduced coordinates and integer
modes while physical radii use $A\mathbf n$.

Agreement between the two implementations checks numerical reconstruction but does not
turn a finite mesh into a continuum result or a synthetic parent into a material model.
The geometry-only, frozen-parent, and radial-response calculations are therefore
interpreted within their declared finite meshes, shell caps, and response laws.
